\documentclass[11pt]{article}
\usepackage[margin=1in]{geometry}
\usepackage{xcolor}
\usepackage{tcolorbox}
\usepackage{hyperref}
\usepackage{enumitem}
\setlist{noitemsep,topsep=0pt}
\usepackage{booktabs}
\usepackage{array}
\usepackage{amsmath}

% Define OSU orange and AI blue colors
\definecolor{osaorange}{RGB}{220,115,38}
\definecolor{aiblue}{RGB}{30,144,255}

% Configure hyperref colors
\hypersetup{
    colorlinks=true,
    linkcolor=blue,
    urlcolor=blue,
    citecolor=blue
}

\title{}
\author{}
\date{}

\begin{document}

% Orange header box with black outline
\begin{tcolorbox}[colback=osaorange,colframe=black,boxrule=2pt,arc=0pt,left=6pt,right=6pt,top=10pt,bottom=10pt]
\centering
\textcolor{white}{\textbf{\Large Week 6 In-Class Worksheet}}\\
\textcolor{white}{\textbf{\large Power, Two-Sample Tests, and ANOVA}}
\end{tcolorbox}

\vspace{8pt}

\noindent\textbf{Name:} \rule{8cm}{0.4pt} \hfill \textbf{Date:} \rule{4cm}{0.4pt}\\[0.5ex]
\textbf{Course:} H524 Introduction to Biostatistics\\
\textbf{Topic:} Statistical Power, Two-Sample Comparisons, and Analysis of Variance

\vspace{8pt}

\section*{Instructions}
Work with a partner to solve the following problems. Show all your work. You may use R to verify your answers and create visualizations, but demonstrate your understanding of the concepts by showing calculations and explaining your reasoning.

\vspace{12pt}

\section{Problem 1: Statistical Power}

A pharmaceutical company is planning a clinical trial to test if a new drug reduces systolic blood pressure. Based on previous studies:
\begin{itemize}
\item Population mean blood pressure: 140 mmHg
\item Standard deviation: 18 mmHg
\item They want to detect a reduction to 135 mmHg (5 mmHg decrease)
\item Significance level: $\alpha = 0.05$ (two-sided test)
\end{itemize}

\vspace{6pt}

\textbf{Part A:} If they enroll 50 patients, what is the power of their study? Use R to calculate.

\vspace{0.3cm}

\textbf{R code:}

\vspace{2.5cm}

\textbf{Power:} \rule{3cm}{0.4pt}

\vspace{0.3cm}

\textbf{Part B:} What sample size is needed to achieve 80\% power?

\vspace{0.3cm}

\textbf{R code:}

\vspace{2.5cm}

\textbf{Required n:} \rule{3cm}{0.4pt}

\vspace{0.3cm}

\textbf{Part C:} What does ``80\% power'' mean in the context of this study? Write a complete sentence.

\vspace{2.5cm}

\textbf{Part D:} If the company can only afford to enroll 100 patients, what is the smallest effect (reduction in mmHg) they could detect with 80\% power?

\vspace{0.3cm}

Hint: Set \texttt{n = 100}, \texttt{power = 0.80}, and solve for \texttt{delta}.

\vspace{2.5cm}

\newpage

\section{Problem 2: Two-Sample t-Test}

A nutritionist compares the effectiveness of two weight-loss diets. She randomly assigns 12 participants to Diet A and 12 to Diet B. After 8 weeks, weight loss (in kg) is measured:

\vspace{0.3cm}

\begin{center}
\begin{tabular}{ll}
\textbf{Diet A:} & 3.2, 4.5, 2.8, 5.1, 3.9, 4.2, 3.6, 4.8, 3.3, 4.0, 3.7, 4.4 \\
\textbf{Diet B:} & 5.2, 6.1, 4.8, 5.9, 5.5, 6.3, 5.0, 5.7, 5.4, 6.0, 5.3, 5.8
\end{tabular}
\end{center}

\vspace{6pt}

\textbf{Part A:} Calculate summary statistics for each group. Use R.

\vspace{0.3cm}

\textbf{Diet A:} $\bar{x}_A = $ \rule{2cm}{0.4pt}, $s_A = $ \rule{2cm}{0.4pt}, $n_A = $ \rule{1.5cm}{0.4pt}

\vspace{0.3cm}

\textbf{Diet B:} $\bar{x}_B = $ \rule{2cm}{0.4pt}, $s_B = $ \rule{2cm}{0.4pt}, $n_B = $ \rule{1.5cm}{0.4pt}

\vspace{0.3cm}

\textbf{Part B:} State the null and alternative hypotheses for testing if Diet B results in greater weight loss than Diet A.

\vspace{0.3cm}

$H_0$: \rule{8cm}{0.4pt}

\vspace{0.5cm}

$H_A$: \rule{8cm}{0.4pt}

\vspace{0.5cm}

\textbf{Part C:} Perform a two-sample t-test in R. Use \texttt{var.equal = TRUE} (assume equal variances). Use a one-sided test.

\vspace{0.3cm}

\textbf{R code:}

\vspace{3cm}

\textbf{Test statistic:} $t = $ \rule{2cm}{0.4pt}

\vspace{0.3cm}

\textbf{Degrees of freedom:} $df = $ \rule{2cm}{0.4pt}

\vspace{0.3cm}

\textbf{p-value:} \rule{2cm}{0.4pt}

\vspace{0.3cm}

\textbf{Part D:} What is your conclusion at $\alpha = 0.05$? Write a complete sentence in context.

\vspace{3cm}

\textbf{Part E:} Calculate and interpret a 95\% confidence interval for the difference in mean weight loss ($\mu_B - \mu_A$).

\vspace{0.3cm}

\textbf{95\% CI:} \rule{4cm}{0.4pt}

\vspace{0.3cm}

\textbf{Interpretation:}

\vspace{2.5cm}

\newpage

\section{Problem 3: Paired t-Test}

A physical therapist measures knee flexibility (degrees of extension) in 10 patients before and after a 6-week treatment program:

\vspace{0.3cm}

\begin{center}
\begin{tabular}{ccc}
\toprule
Patient & Before & After \\
\midrule
1 & 115 & 125 \\
2 & 110 & 120 \\
3 & 120 & 128 \\
4 & 105 & 115 \\
5 & 118 & 130 \\
6 & 112 & 122 \\
7 & 108 & 118 \\
8 & 122 & 135 \\
9 & 114 & 124 \\
10 & 116 & 126 \\
\bottomrule
\end{tabular}
\end{center}

\vspace{6pt}

\textbf{Part A:} Calculate the differences (After - Before) for each patient.

\vspace{0.3cm}

Differences: \rule{10cm}{0.4pt}

\vspace{0.5cm}

\textbf{Part B:} Calculate the mean and standard deviation of the differences.

\vspace{0.3cm}

$\bar{d} = $ \rule{2cm}{0.4pt} \hspace{2cm} $s_d = $ \rule{2cm}{0.4pt}

\vspace{0.5cm}

\textbf{Part C:} State the hypotheses for testing if flexibility improves after treatment.

\vspace{0.3cm}

$H_0$: \rule{8cm}{0.4pt}

\vspace{0.5cm}

$H_A$: \rule{8cm}{0.4pt}

\vspace{0.5cm}

\textbf{Part D:} Perform a paired t-test in R.

\vspace{0.3cm}

\textbf{R code:}

\vspace{3cm}

\textbf{Test statistic:} $t = $ \rule{2cm}{0.4pt} \hspace{1cm} \textbf{p-value:} \rule{2cm}{0.4pt}

\vspace{0.5cm}

\textbf{Part E:} What is your conclusion? Does the treatment significantly improve flexibility?

\vspace{2.5cm}

\textbf{Part F:} Why would it be wrong to use an independent two-sample t-test for this data? Explain.

\vspace{3cm}

\newpage

\section{Problem 4: One-Way ANOVA}

A researcher tests the effect of four different exercise programs on weight loss (kg) over 3 months. She randomly assigns participants to one of four programs:

\vspace{0.3cm}

\begin{center}
\begin{tabular}{cccc}
\toprule
Walking & Swimming & Cycling & Gym \\
\midrule
3.2 & 4.5 & 5.8 & 6.2 \\
2.8 & 4.8 & 6.1 & 6.5 \\
3.5 & 4.2 & 5.5 & 5.8 \\
3.0 & 4.6 & 6.0 & 6.3 \\
3.3 & 4.4 & 5.7 & 6.0 \\
\midrule
Mean: 3.16 & Mean: 4.50 & Mean: 5.82 & Mean: 6.16 \\
SD: 0.27 & SD: 0.23 & SD: 0.24 & SD: 0.27 \\
\bottomrule
\end{tabular}
\end{center}

\vspace{6pt}

\textbf{Part A:} State the null and alternative hypotheses for an ANOVA test.

\vspace{0.3cm}

$H_0$: \rule{10cm}{0.4pt}

\vspace{0.5cm}

$H_A$: \rule{10cm}{0.4pt}

\vspace{0.5cm}

\textbf{Part B:} Why is it better to use ANOVA instead of multiple pairwise t-tests?

\vspace{3cm}

\textbf{Part C:} Perform a one-way ANOVA in R.

\vspace{0.3cm}

\textbf{R code to create data frame and run ANOVA:}

\vspace{4cm}

\textbf{F-statistic:} $F = $ \rule{2cm}{0.4pt}

\vspace{0.3cm}

\textbf{Degrees of freedom:} $df_1 = $ \rule{1.5cm}{0.4pt}, $df_2 = $ \rule{1.5cm}{0.4pt}

\vspace{0.3cm}

\textbf{p-value:} \rule{2cm}{0.4pt}

\vspace{0.5cm}

\textbf{Part D:} What is your conclusion from the ANOVA? Does it appear that exercise programs differ in effectiveness?

\vspace{2.5cm}

\textbf{Part E:} If the ANOVA is significant, perform Tukey's HSD post-hoc test to determine which programs differ from each other.

\vspace{0.3cm}

\textbf{R code:}

\vspace{2cm}

\textbf{Which pairwise comparisons are significant?} (List them)

\vspace{2.5cm}

\textbf{Part F:} Based on your analysis, which exercise program would you recommend for weight loss? Justify your answer.

\vspace{3cm}

\newpage

\section{Problem 5: Choosing the Right Test}

For each scenario below, identify which statistical test is most appropriate and explain why.

\vspace{0.5cm}

\textbf{Scenario A:} A doctor measures blood glucose in 15 diabetic patients before and after taking a new medication.

\vspace{0.3cm}

\textbf{Test:} \rule{6cm}{0.4pt}

\vspace{0.3cm}

\textbf{Reason:}

\vspace{2cm}

\textbf{Scenario B:} A researcher compares average IQ scores between children who attended preschool (n=25) and children who did not attend preschool (n=30).

\vspace{0.3cm}

\textbf{Test:} \rule{6cm}{0.4pt}

\vspace{0.3cm}

\textbf{Reason:}

\vspace{2cm}

\textbf{Scenario C:} A study compares patient satisfaction scores across five different hospitals.

\vspace{0.3cm}

\textbf{Test:} \rule{6cm}{0.4pt}

\vspace{0.3cm}

\textbf{Reason:}

\vspace{2cm}

\textbf{Scenario D:} An agronomist compares crop yield for two fertilizer brands, using 20 fields. Each field is split in half, with one fertilizer applied to each half.

\vspace{0.3cm}

\textbf{Test:} \rule{6cm}{0.4pt}

\vspace{0.3cm}

\textbf{Reason:}

\vspace{2cm}

\textbf{Scenario E:} A pharmaceutical company wants to determine how many patients to enroll in a trial to detect a 15-point reduction in pain score with 85\% power.

\vspace{0.3cm}

\textbf{Test/Analysis:} \rule{6cm}{0.4pt}

\vspace{0.3cm}

\textbf{Reason:}

\vspace{2cm}

\newpage

\section*{Reflection Questions}

\textbf{1.} Explain the relationship between Type I error, Type II error, and power. How are they connected?

\vspace{3cm}

\textbf{2.} What are the three main factors that affect statistical power, and how does each one influence it?

\vspace{3cm}

\textbf{3.} When should you use a paired t-test instead of an independent two-sample t-test? Give an example.

\vspace{3cm}

\textbf{4.} What information does an ANOVA F-test give you? What doesn't it tell you?

\vspace{3cm}

\textbf{5.} Why is it important to calculate power and sample size \textit{before} conducting a study rather than after?

\vspace{3cm}

\end{document}
